\section{Training Methods and Systems}
\sectionkicker{TRAINING}
\begin{wideflow}
{\Large\bfseries 鍛えの手筋――ContextLMの式とOlmo-core 3の仕組み}\par
\smallskip
{\color{SurveyMuted}編集できる綴じと群れの技を、式と分母と土台つきで読む。今号いちばんの読みどころである。}\par
\end{wideflow}

ContextLMは、文脈を綴じとして編む\cite{contextlm}。Bashで擦り合わせた綴じの品で、使い回す文頭の再利用をフロップスという物差しで量る。式3の置き場は「最初の食い違いより後の文の詰め直しと、新しく生む字句の解きの勘定」であり、壁時計でも当たり率でもない。一致した文頭は勘定に入らないという約束が要である。ここを外すと後の割り引きの話が全部ずれる。

式5は技の進め方の記述である。重みは固めたまま、綴じの文書 \(s\) を鍛え、鍛えと試しと取って置きの三分けで運び、取って置きでは一度だけ量る。持ち越しの伸びは最大で35.9点と置かれる。式の本体
\[ s^* = \arg\max_s E_{x\sim D}[R(\tau(x;s))] \]
は、文脈の中に技を置く流儀の最適化であり、重みを学ぶ流儀ではない。式6は別物で、うまくいった軌跡が二本に満たない群では効率の項は黙り、成否の項だけが残る。成功者の中を安い順に並べ替える仕掛けであり、正誤の裁きは別が担う。式4は確かめたとおりで動かない。式5の全文引用の範囲を超える細目は言わない。

数の置き場は混ぜない。BCPは59.4\%で相対11.4\%の伸び、向かうはCodex流の要約という最強の土台である。使い回しの割り引きは同土台に21.5\%、MEM1に28.9\%で、両者は統合しない。TB2.1は同等の当たりで7割の勘定、TBLiteは73.7対67.0で9割1分の勘定で、土台は同じCodex流である。EdgeBenchの12時間10題は5\%の伸びに59\%の割り引き、群れは65\%の伸び、オンラインの強化学習はQwen3.5-9Bで28.8から42.5へと13.7差、同じ手順のCodex harnessには0.4点と38.8\%の割り引きで並ぶ。いずれも紙の言い分であり、許諾は紙がBY-4.0、置き場がBY-NC-4.0である。

Olmo-core 3は、群れをどう捌くかの仕組みである\cite{olmocore3}。DDPに技持ちを住まわせ、GPUの上で回し先を決めて並べ替えの余分を抑え、束ねGEMMとMXFP8で固める。MXFP8の効きはB300四枚の均一負荷という構えでの話で、通りが21\%伸び、103GiBが95GiBに痩せると置かれる。他の構えへの広げはしない。技持ちを8から128に増やし、一手に4つ動かして3.2Bほど働かせても、通りは5\%も落ちない。8枚B300の予備走りでは旧実装の2.7倍ほどと置かれる。

物差しの置き場は三つに分ける。47Bは測った値、1.2Tは回し（能動58.36B・512枚・最大858TFLOP/s/GPU、でたらめ回しの条件の仕組みの値であり、品の良し悪しではない）、2.38Tは短い容量試しである。三つを一つの大きさに束ねてはならない。うまくいかなかった話も載せる。回しの釣り合いの点は失敗に終わり、量り方への戒めとして残る。報告書の本体は手が届かず、紙の付録や絵の言及のない中身は言わない。

\begin{claimboundary}[資料と評価の境界]
式の読みは一次資料の範囲であり、全文引用の範囲を超える細目は言わない。数の土台は混ぜず、相対と差分、紙と置き場の許諾を切り分ける。三つの大きさは別の置き場であり、一つに束ねない。報告書の本体や付録の言及のない中身は言わない。
\end{claimboundary}
